\documentclass{ximera}

\input{../../preamble.tex}

\title{Angles}

\begin{document}

\begin{abstract}
inverses
\end{abstract}
\maketitle




We have two types of trigonometric functions.


We have functions that map angles to ratios:


\[
\sin(\theta), \cos(\theta) :  \text{ Angles } \to \text{ Ratios }
\]


And, we have their inverses:



\[
\sin^{-1}(t), \cos^{-1}(t) :  \text{ Ratios } \to \text{ Angles }
\]








Your calculator has buttons for all of these.


\textbf{Note:}  Your calculator has two modes for angle measurements: degrees and radians.  Remember to switch between these modes when calculating.



\begin{question} 

Approximate the following.

\begin{itemize}
\item $\cos(20^{\circ}) \approx \answer[tolerance=0.001]{0.9396926208}$
\item $\sin(65^{\circ}) \approx \answer[tolerance=0.001]{0.906307787}$
\item $\cos\left( \frac{\pi}{3} \right) \approx \answer[tolerance=0.001]{0.5}$
\item $\sin\left( \frac{\pi}{4} \right) \approx \answer[tolerance=0.001]{0.7071067812}$
\end{itemize}

\end{question}





Your calculator also has buttons for the inverse functions. 




You can also supply the calculator with the ratio and the calculator returns the angle.



\begin{question} 

Approximate the following.

\begin{itemize}
\item $\cos^{-1}(0.9396926208) \approx \answer[tolerance=0.001]{20}^{\circ}$
\item $\sin^{-1}(0.906307787) \approx \answer[tolerance=0.001]{65}^{\circ}$
\item $\cos^{-1}(0.5) \approx \answer[tolerance=0.001]{1.047197551}$
\item $\sin^{-1}(0.7071067812) \approx \answer[tolerance=0.001]{0.7853981634}$
\end{itemize}

\end{question}




\subsection*{Other Quadrants}

Remember that the signs of sine and cosine changes as we move through the four quadrants.  We introduced a reference angle and a refrence triangle to help with lengths and then we also had to think of the sign. 

The calculator also knows that the signs change. 

The calculator will give the correct values of sine and cosine along with their proper sign.  However, going backwards is a different story. 

If you give the calculator a value of sine and ask it for the angle, then there are an infinite number of angles that have that same value of sine.  Just keep spinning around the circle. 



\begin{example} Going Backwards with Sine

$\blacktriangleright$ Given  $\sin(\theta) = \frac{1}{2}$, then any of the following are possible.


\begin{itemize}
\item $\theta = -\frac{7\pi}{6}$
\item $\theta = \frac{\pi}{6}$
\item $\theta = \frac{5\pi}{6}$
\item $\theta = \frac{13\pi}{6}$
\item $\theta = \frac{17\pi}{6}$
\item an infinite number of other angles
\end{itemize}


\end{example}



Therefore, the calculator just focuses on two quadrants, one for each sign. 


For sine, the calculator gives angles in the fourth and first quadrants in the interval $\left( -\frac{\pi}{2}, \frac{\pi}{2}  \right)$. 

\begin{itemize}
\item When you supply a negative value for sine, the calculator returns an angle in the interval $\left( -\frac{\pi}{2}, 0 \right)$, which is the fourth quadrant.
\item When you supply a positive value for sine, the calculator returns an angle in the interval $\left( 0, \frac{\pi}{2}, 0 \right)$, which is the first quadrant.
\end{itemize}


If the angle you desire is actually in another quadrant, then it is up to you to figure out its actual value from there.







\begin{question} 

$150^{\circ}$ is an angle in the second quadrant. $\sin(150^{\circ}) = \frac{1}{2} = 0.5$.


Now, ask your calculator to give you an angle whose sine is $0.5$.


$\sin^{-1}(0.5) = \answer{30}^{\circ}$


It doesn't give $150^{\circ}$.  It gives $30^{\circ}$.

\end{question}














\begin{example} Going Backwards with Cosine

$\blacktriangleright$ Given  $\cos(\theta) = \frac{1}{2}$, then any of the following are possible.


\begin{itemize}
\item $\theta = \frac{\pi}{3}$
\item $\theta = -\frac{\pi}{3}$
\item $\theta = \frac{13\pi}{6}$
\item $\theta = -\frac{13\pi}{6}$
\item $\theta = \frac{25\pi}{6}$
\item an infinite number of other angles
\end{itemize}


\end{example}



Therefore, the calculator just focuses on two quadrants, one for each sign. 


For cosine, the calculator gives angles in the first and second quadrants in the interval $(0, \pi)$. 

\begin{itemize}
\item When you supply a positive value for cosine, the calculator returns an angle in the interval $\left( 0, \frac{\pi}{2} \right)$, which is the first quadrant.
\item When you supply a negative value for cosine, the calculator returns an angle in the interval $\left( \frac{\pi}{2}, \pi \right)$, which is the second quadrant.
\end{itemize}


If the angle you desire is actually in another quadrant, then it is up to you to figure out its actual value from there.







\begin{question} 

$240^{\circ}$ is an angle in the third quadrant. $\cos(240^{\circ}) = -\frac{1}{2} = 0.5$.


Now, ask your calculator to give you an angle whose cosine is $-0.5$.


$\cos^{-1}(-0.5) = \answer{120}^{\circ}$


It doesn't give $240^{\circ}$.  It gives $120^{\circ}$.

\end{question}














We'll study this more later in the course.  The short story is we have trigonometric functions and their inverse functions.




\begin{itemize}
\item $\sin^{-1}(x)$ will give angles in the fourth and first quadrants.
\item $\cos^{-1}(x)$ will give angles in the first and second quadrants.
\item $\tan^{-1}(x)$ will give angles in the fourth and first quadrants.
\end{itemize}


The above is certainly proper inverse function notation, but it confuses everyone, because it is also proper notation for a reciprocal.  Therefore, we prefer other names:




\begin{itemize}
\item \textbf{\textcolor{blue!55!black}{$\arcsin(x)$}}  will give angles in the fourth and first quadrants.
\item \textbf{\textcolor{blue!55!black}{$\arccos(x)$}}  will give angles in the first and second quadrants.
\item \textbf{\textcolor{blue!55!black}{$\arctan(x)$}}  will give angles in the fourth and first quadrants.
\end{itemize}





\begin{center}
\textbf{\textcolor{green!50!black}{ooooo-=-=-=-ooOoo-=-=-=-ooooo}} \\

more examples can be found by following this link\\ \link[More Examples of Right Triangles]{https://ximera.osu.edu/csccmathematics/precalculus2/precalculus2/rightTriangles/examples/exampleList}

\end{center}






\end{document}
